%2multibyte Version: 5.50.0.2960 CodePage: 65001
% !TeX root = ./Slides_Lez_10_Python_Rischio_Credito.tex
% ------------------------------------------------------------
% Macro matematiche ufficiali del progetto MQF
% Coerenti con MQF_Notazione.tex
% ------------------------------------------------------------
% Fallback utile per compilazione esterna se qualche titolo resta scritto come \QTR.
% Scientific WorkPlace usa \QTR internamente; in Beamer \QTR{frametitle}{Titolo}
% deve espandersi in \frametitle{Titolo}.

\documentclass[handout]{beamer}%
\usepackage{mathpazo}
%\usepackage{hyperref}
\usepackage{multimedia}
\usepackage{amsmath}
\usepackage{amsfonts}
\usepackage{amssymb}
\usepackage{graphicx}%
\setcounter{MaxMatrixCols}{30}
%TCIDATA{OutputFilter=latex2.dll}
%TCIDATA{Version=5.50.0.2960}
%TCIDATA{Codepage=65001}
%TCIDATA{CSTFile=beamer.cst}
%TCIDATA{Created=Saturday, May 16, 2026 09:09:59}
%TCIDATA{LastRevised=Saturday, May 16, 2026 11:24:58}
%TCIDATA{<META NAME="GraphicsSave" CONTENT="32">}
%TCIDATA{<META NAME="SaveForMode" CONTENT="2">}
%TCIDATA{BibliographyScheme=Manual}
%TCIDATA{<META NAME="DocumentShell" CONTENT="Other Documents\SW\Slides - Beamer">}
%BeginMSIPreambleData
\graphicspath{{../graphics/}}
\providecommand{\U}[1]{\protect\rule{.1in}{.1in}}
%EndMSIPreambleData
\newenvironment{stepenumerate}{\begin{enumerate}[<+->]}{\end{enumerate}}
\newenvironment{stepitemize}{\begin{itemize}[<+->]}{\end{itemize} }
\newenvironment{stepenumeratewithalert}{\begin{enumerate}[<+-| alert@+>]}{\end{enumerate}}
\newenvironment{stepitemizewithalert}{\begin{itemize}[<+-| alert@+>]}{\end{itemize} }
\usetheme{Madrid}


\setbeamertemplate{footline}
{
  \leavevmode%
  \hbox{%
    \begin{beamercolorbox}[wd=.20\paperwidth,ht=2.25ex,dp=1ex,center]
      {author in head/foot}%
      \usebeamerfont{author in head/foot}\insertshortauthor
    \end{beamercolorbox}%
    \begin{beamercolorbox}[wd=.60\paperwidth,ht=2.25ex,dp=1ex,center]
      {title in head/foot}%
      \usebeamerfont{title in head/foot}\insertshorttitle
    \end{beamercolorbox}%
    \begin{beamercolorbox}[wd=.20\paperwidth,ht=2.25ex,dp=1ex,right]
      {date in head/foot}%
      \usebeamerfont{date in head/foot}%
      \insertshortdate{}\hspace*{2em}%
      \insertframenumber{} / \inserttotalframenumber\hspace*{2ex}
    \end{beamercolorbox}%
  }%
  \vskip0pt%
}




\providecommand{\R}{\mathbb{R}}
\providecommand{\N}{\mathbb{N}}
\providecommand{\E}{\mathbb{E}}
\providecommand{\Prob}{\mathbb{P}}
\providecommand{\Q}{\mathbb{Q}}
\providecommand{\F}{\mathcal{F}}
\providecommand{\B}{\mathcal{B}}
\providecommand{\G}{\mathcal{G}}
\providecommand{\1}{\mathbf{1}}
\providecommand{\Var}{\operatorname{Var}}
\providecommand{\Cov}{\operatorname{Cov}}
\providecommand{\VaR}{\operatorname{VaR}}
\providecommand{\CVaR}{\operatorname{CVaR}}
\providecommand{\argmin}{\operatorname*{arg\,min}}
\providecommand{\argmax}{\operatorname*{arg\,max}}
\providecommand{\EVPI}{\operatorname{EVPI}}
\providecommand{\VSS}{\operatorname{VSS}}
\providecommand{\QTR}[2]{\csname #1\endcsname{#2}}
%BeginMSIPreambleData
\ifx\pdfoutput\relax\let\pdfoutput=\undefined\fi
\newcount\msipdfoutput
\ifx\pdfoutput\undefined\else
\ifcase\pdfoutput\else
\msipdfoutput=1
\ifx\paperwidth\undefined\else
\ifdim\paperheight=0pt\relax\else\pdfpageheight\paperheight\fi
\ifdim\paperwidth=0pt\relax\else\pdfpagewidth\paperwidth\fi
\fi\fi\fi
%EndMSIPreambleData

%******************
% Slide 1
%******************

\title[Lez. 10 - Python, rischio di credito]{Lezione 10 - Python, rischio di credito}
\subtitle{Metodi Quantitativi per la Finanza}
\author[P. Falbo, C. Pelizzari]{{\LARGE Paolo Falbo, Cristian Pelizzari}}
\institute[]{{\large Dipartimento di Economia e Management}}
\date[]{}
\titlegraphic{
\includegraphics[width=3.3cm]{../graphics/logo-Unibs-esteso.png}
}

\begin{document}
\maketitle



\section{Apertura della lezione}

%******************
% Slide 2
%******************

%TCIMACRO{\TeXButton{BeginFrame}{\begin{frame}}}%
%BeginExpansion
\begin{frame}%
%EndExpansion

\QTR{frametitle}{Obiettivi della lezione}

%TCIMACRO{\TeXButton{Transparency}{\setbeamercovered{transparent=20}}}%
%BeginExpansion
\setbeamercovered{transparent=20}%
%EndExpansion

La Lezione 10 ha tre obiettivi irrinunciabili.

\begin{enumerate}

\item \textbf{Un piccolo approfondimento Python}

Utilizzare maschere booleane e selezioni condizionali per rendere
operativi eventi definiti probabilisticamente.

\item \textbf{Il Caso Aula: China Evergrande 2021}

Trasformare la teoria delle catene di Markov e delle misure di rischio
in un modello Monte Carlo per la perdita di un portafoglio creditizio.

\item \textbf{Il Caso TakeHome: Lehman Brothers 2008}

Trasferire autonomamente la stessa architettura metodologica a un
problema finanziario distinto, che non verr\`a sviluppato in aula.

\end{enumerate}

\vspace{0.3cm}

Il protocollo generale delle lezioni applicative, introdotto nella
Lezione 4, verr\`a richiamato solo quando necessario.

%TCIMACRO{\TeXButton{Transition: Box Out}{\transboxout}}}%
%BeginExpansion
\transboxout%
%EndExpansion
%TCIMACRO{\TeXButton{EndFrame}{\end{frame}}}%
%BeginExpansion
\end{frame}%
%EndExpansion


%******************
% Slide 3
%******************

%TCIMACRO{\TeXButton{BeginFrame}{\begin{frame}}}%
%BeginExpansion
\begin{frame}%
%EndExpansion

\QTR{frametitle}{Dove siamo nel corso}

%TCIMACRO{\TeXButton{Transparency}{\setbeamercovered{transparent=20}}}%
%BeginExpansion
\setbeamercovered{transparent=20}%
%EndExpansion

La Lezione 10 porta sul piano computazionale gli strumenti sviluppati
nei due capitoli precedenti.

\vspace{0.3cm}

\begin{description}

\item[\textbf{Capitolo 8}]
catene di Markov, matrici di transizione, distribuzioni degli stati e
dinamica multi-periodale.

\vspace{0.2cm}

\item[\textbf{Capitolo 9}]
migrazioni creditizie, distribuzione delle perdite e misure di rischio,
con particolare attenzione a VaR e CVaR.

\end{description}

\vspace{0.4cm}

Nella lezione di oggi, il passaggio centrale \`e:

\[
\text{regime sistemico}
\longrightarrow
\text{migrazioni creditizie}
\longrightarrow
\text{perdite individuali}
\]
\vspace{-0.2cm}
\[
\longrightarrow
\text{perdita di portafoglio}
\longrightarrow
\text{distribuzione simulata}
\longrightarrow
\text{VaR e CVaR}.
\]



%\begin{center}
%\textbf{Obiettivo: rendere simulabile e verificabile il collegamento
%tra dinamica creditizia e rischio di portafoglio.}
%\end{center}

%TCIMACRO{\TeXButton{Transition: Box Out}{\transboxout}}}%
%BeginExpansion
\transboxout%
%EndExpansion
%TCIMACRO{\TeXButton{EndFrame}{\end{frame}}}%
%BeginExpansion
\end{frame}%
%EndExpansion


%============================================================
% SEZIONE 1 — PICCOLO FOCUS PYTHON
%============================================================

\section{Piccolo focus Python}

%******************
% Slide 4
%******************

%TCIMACRO{\TeXButton{BeginFrame}{\begin{frame}}}%
%BeginExpansion
\begin{frame}%
%EndExpansion

\QTR{frametitle}{Da un evento probabilistico a una condizione logica}

%TCIMACRO{\TeXButton{Transparency}{\setbeamercovered{transparent=20}}}%
%BeginExpansion
\setbeamercovered{transparent=20}%
%EndExpansion

Nel Caso Aula, considereremo, tra gli altri, l'evento

\[
\{\exists t : M_t=C\},
\]

cio\`e l'evento in cui la traiettoria del regime sistemico visita almeno
una volta lo stato di crisi.

\vspace{0.4cm}

In una simulazione Monte Carlo, disponiamo di molte traiettorie:

\[
M^{(1)},M^{(2)},\ldots,M^{(N_{\mathrm{MC}})}.
\]

Per ciascuna replica dobbiamo quindi rispondere a una domanda logica:

\begin{center}
\textbf{la traiettoria contiene almeno una volta lo stato \(C\)?}
\end{center}

\vspace{0.3cm}

La risposta \`e: \texttt{True} \text{oppure} \texttt{False}.

\vspace{0.2cm}

Python traduce cos\`i un evento probabilistico in una condizione
verificabile replica per replica.

%TCIMACRO{\TeXButton{Transition: Box Out}{\transboxout}}}%
%BeginExpansion
\transboxout%
%EndExpansion
%TCIMACRO{\TeXButton{EndFrame}{\end{frame}}}%
%BeginExpansion
\end{frame}%
%EndExpansion


%******************
% Slide 5
%******************

%TCIMACRO{\TeXButton{BeginFrame}{\begin{frame}}}%
%BeginExpansion
\begin{frame}%
%EndExpansion

\QTR{frametitle}{Maschere booleane e \texttt{np.any}}

%TCIMACRO{\TeXButton{Transparency}{\setbeamercovered{transparent=20}}}%
%BeginExpansion
\setbeamercovered{transparent=20}%
%EndExpansion

Supponiamo che l'array \texttt{M} contenga le traiettorie simulate del
regime sistemico:

\[
\texttt{M.shape}
=
(N_{\mathrm{MC}},T).
\]

Il confronto \texttt{M == C} produce un array booleano della stessa dimensione.

\vspace{0.3cm}

Per verificare se una traiettoria visita almeno una volta lo stato \(C\)
possiamo usare:

\begin{center}
\texttt{crisis = np.any(M == C, axis=1)}
\end{center}

\vspace{0.3cm}

Il parametro \texttt{axis=1} indica che il controllo viene effettuato
lungo il tempo, separatamente per ciascuna replica.

\vspace{0.3cm}

L'output \`e un vettore:

\[
\texttt{crisis}
=
(\texttt{True},\texttt{False},\ldots),
\]

con una componente per ogni scenario Monte Carlo.

%TCIMACRO{\TeXButton{Transition: Box Out}{\transboxout}}}%
%BeginExpansion
\transboxout%
%EndExpansion
%TCIMACRO{\TeXButton{EndFrame}{\end{frame}}}%
%BeginExpansion
\end{frame}%
%EndExpansion


%******************
% Slide 6
%******************

%TCIMACRO{\TeXButton{BeginFrame}{\begin{frame}}}%
%BeginExpansion
\begin{frame}%
%EndExpansion

\QTR{frametitle}{Selezionare sottocampioni con una maschera}

%TCIMACRO{\TeXButton{Transparency}{\setbeamercovered{transparent=20}}}%
%BeginExpansion
\setbeamercovered{transparent=20}%
%EndExpansion

Una maschera booleana pu\`o essere utilizzata per selezionare solo le
osservazioni che soddisfano una determinata condizione.



Se \texttt{loss} contiene le perdite di portafoglio simulate, allora:

\begin{center}
\texttt{loss[crisis]}
\end{center}

seleziona le perdite associate agli scenari in cui il regime sistemico
entra almeno una volta nello stato \(C\).

Analogamente:

\begin{center}
\texttt{loss[{\raise.17ex\hbox{$\scriptstyle\sim$}}crisis]}
\end{center}

seleziona gli scenari senza ingresso nello stato di crisi.



Sul piano probabilistico stiamo costruendo due sottocampioni:

\[
L^P \mid \text{Crisi}
\qquad\text{e}\qquad
L^P \mid \text{No Crisi}.
\]

\vspace{0.3cm}

\begin{center}
\textbf{La maschera booleana rende operativa una partizione del campione
definita dal modello probabilistico.}
\end{center}

%TCIMACRO{\TeXButton{Transition: Box Out}{\transboxout}}}%
%BeginExpansion
\transboxout%
%EndExpansion
%TCIMACRO{\TeXButton{EndFrame}{\end{frame}}}%
%BeginExpansion
\end{frame}%
%EndExpansion


%============================================================
% SEZIONE 2 — DAI CAPITOLI 8–9 AL PROBLEMA COMPUTAZIONALE
%============================================================

\section{Dai Capitoli 8--9 al problema computazionale}

%******************
% Slide 7
%******************

%TCIMACRO{\TeXButton{BeginFrame}{\begin{frame}}}%
%BeginExpansion
\begin{frame}%
%EndExpansion

\QTR{frametitle}{Dal Capitolo 8: il regime sistemico come catena di Markov}

%TCIMACRO{\TeXButton{Transparency}{\setbeamercovered{transparent=20}}}%
%BeginExpansion
\setbeamercovered{transparent=20}%
%EndExpansion

Nel Caso Aula, il contesto macro-finanziario viene rappresentato da un
processo a stati discreti

\[
M_t \in \{N,S,C\},
\]

dove:

\begin{itemize}

\item \(N\) indica un regime normale;

\item \(S\) un regime di stress;

\item \(C\) un regime di crisi.

\end{itemize}

%\vspace{0.1cm}

L'evoluzione del regime \`e descritta da una matrice di transizione

\[
Q=(q_{rs}),
\]

con stato iniziale $M_0=S$.

%\vspace{0.3cm}

Il primo problema computazionale consiste quindi nel simulare
molte traiettorie $M_0,M_1,\ldots,M_T$.

\begin{center}
\textbf{Il Capitolo 8 fornisce il modello; Python genera le traiettorie.}
\end{center}

%TCIMACRO{\TeXButton{Transition: Box Out}{\transboxout}}}%
%BeginExpansion
\transboxout%
%EndExpansion
%TCIMACRO{\TeXButton{EndFrame}{\end{frame}}}%
%BeginExpansion
\end{frame}%
%EndExpansion


%******************
% Slide 8
%******************

%TCIMACRO{\TeXButton{BeginFrame}{\begin{frame}}}%
%BeginExpansion
\begin{frame}%
%EndExpansion

\QTR{frametitle}{Dal Capitolo 8: migrazioni creditizie condizionate al regime}

%TCIMACRO{\TeXButton{Transparency}{\setbeamercovered{transparent=20}}}%
%BeginExpansion
\setbeamercovered{transparent=20}%
%EndExpansion

Per ciascun debitore \(i\), il rating evolve secondo un processo $X_{i,t}$.

Nel Caso Aula, gli stati creditizi sono:

\[
\{BBB,BB,B,CCC,D\},
\]

con \(D\) stato assorbente.

%\vspace{0.3cm}

La matrice di transizione dipende dal regime sistemico corrente:

\[
P^{(N)},\qquad
P^{(S)},\qquad
P^{(C)}.
\]

In particolare,

\[
\Pr
\left(
X_{i,t+1}=j
\mid
X_{i,t}=k,M_t=m
\right)
=
p_{kj}^{(m)}.
\]

\vspace{0.3cm}

Quindi il regime \(M_t\) determina quale matrice di migrazione
utilizzare nella transizione successiva.

\begin{center}
\textbf{La dinamica creditizia individuale \`e condizionata dallo
stato sistemico comune.}
\end{center}

%TCIMACRO{\TeXButton{Transition: Box Out}{\transboxout}}}%
%BeginExpansion
\transboxout%
%EndExpansion
%TCIMACRO{\TeXButton{EndFrame}{\end{frame}}}%
%BeginExpansion
\end{frame}%
%EndExpansion


%******************
% Slide 9
%******************

%TCIMACRO{\TeXButton{BeginFrame}{\begin{frame}}}%
%BeginExpansion
\begin{frame}%
%EndExpansion

\QTR{frametitle}{Dipendenza sistemica: condizionare cambia la struttura}

%TCIMACRO{\TeXButton{Transparency}{\setbeamercovered{transparent=20}}}%
%BeginExpansion
\setbeamercovered{transparent=20}%
%EndExpansion

In ogni replica Monte Carlo, tutti i debitori condividono la stessa traiettoria sistemica:

\vspace{-0.2cm}

\[
M^{(r)}
=
\left(
M_0^{(r)},M_1^{(r)},\ldots,M_T^{(r)}
\right).
\]

%\vspace{0.3cm}

Condizionatamente a questa traiettoria, le migrazioni dei singoli debitori sono simulate indipendentemente.

%\vspace{0.3cm}

Ma, senza condizionare sul regime sistemico, le perdite dei debitori non sono indipendenti: risentono tutte dello stesso fattore comune.

\[
M_t
\longrightarrow
\begin{cases}
X_{1,t+1}\\
X_{2,t+1}\\
\vdots\\
X_{n,t+1}
\end{cases}
\]

\vspace{0.3cm}

\begin{center}
\textbf{Indipendenza condizionata non significa indipendenza
incondizionata.}
\end{center}

Il regime sistemico introduce cos\`i dipendenza tra le perdite
individuali senza imporre traiettorie creditizie identiche.

%TCIMACRO{\TeXButton{Transition: Box Out}{\transboxout}}}%
%BeginExpansion
\transboxout%
%EndExpansion
%TCIMACRO{\TeXButton{EndFrame}{\end{frame}}}%
%BeginExpansion
\end{frame}%
%EndExpansion


%******************
% Slide 10
%******************

%TCIMACRO{\TeXButton{BeginFrame}{\begin{frame}}}%
%BeginExpansion
\begin{frame}%
%EndExpansion

\QTR{frametitle}{Dal Capitolo 9: dalla migrazione alla perdita monetaria}

%TCIMACRO{\TeXButton{Transparency}{\setbeamercovered{transparent=20}}}%
%BeginExpansion
\setbeamercovered{transparent=20}%
%EndExpansion

La simulazione dei rating non \`e ancora il risultato finale.

\vspace{0.3cm}

Per ogni debitore, dobbiamo trasformare la traiettoria creditizia in una perdita monetaria.

\vspace{0.3cm}

Nel Caso Aula, distingueremo due situazioni:

\begin{description}

\item[\textbf{Nessun default}]
la perdita dipende dalla migrazione tra rating iniziale e rating finale;

\vspace{0.2cm}

\item[\textbf{Default}]
la perdita dipende dall'esposizione e dalla severit\`a della perdita
in caso di default.

\end{description}

\vspace{0.4cm}

Il passaggio concettuale \`e quindi:

\[
\text{traiettoria di rating}
\longrightarrow
\text{perdita individuale}
\longrightarrow
\text{perdita di portafoglio}.
\]

%\begin{center}
%\textbf{Gli stati creditizi diventano quantit\`a monetarie.}
%\end{center}

%TCIMACRO{\TeXButton{Transition: Box Out}{\transboxout}}}%
%BeginExpansion
\transboxout%
%EndExpansion
%TCIMACRO{\TeXButton{EndFrame}{\end{frame}}}%
%BeginExpansion
\end{frame}%
%EndExpansion


%******************
% Slide 11
%******************

%TCIMACRO{\TeXButton{BeginFrame}{\begin{frame}}}%
%BeginExpansion
\begin{frame}%
%EndExpansion

\QTR{frametitle}{Dal Capitolo 9: dalla perdita alle misure di rischio}

%TCIMACRO{\TeXButton{Transparency}{\setbeamercovered{transparent=20}}}%
%BeginExpansion
\setbeamercovered{transparent=20}%
%EndExpansion

Ogni replica Monte Carlo produce una perdita complessiva di portafoglio:

\[
L^{P,(r)},
\qquad
r=1,\ldots,N_{\mathrm{MC}}.
\]

L'insieme delle perdite simulate

\[
\left\{
L^{P,(1)},
L^{P,(2)},
\ldots,
L^{P,(N_{\mathrm{MC}})}
\right\}
\]

costituisce il campione empirico della distribuzione di \(L^P\).

%\vspace{0.3cm}

Su questa distribuzione possiamo calcolare:

\begin{itemize}

\item perdita media;

\item varianza e deviazione standard;

\item VaR a diversi livelli di confidenza;

\item CVaR a diversi livelli di confidenza.

\end{itemize}

%\vspace{0.3cm}

Il percorso completo diventa quindi:

\[
\text{Markov}
\longrightarrow
\text{migrazioni}
\longrightarrow
L^P
\longrightarrow
\text{distribuzione empirica}
\longrightarrow
\text{VaR e CVaR}.
\]

%\begin{center}
%\textbf{Il Caso Aula costruir\`a operativamente questa catena.}
%\end{center}

%TCIMACRO{\TeXButton{Transition: Box Out}{\transboxout}}}%
%BeginExpansion
\transboxout%
%EndExpansion
%TCIMACRO{\TeXButton{EndFrame}{\end{frame}}}%
%BeginExpansion
\end{frame}%
%EndExpansion


%============================================================
% SEZIONE 3 — CASO AULA: EVERGRANDE 2021 — PROBLEMA E MODELLO
%============================================================

\section{Caso Aula: Evergrande 2021}

%******************
% Slide 12
%******************

%TCIMACRO{\TeXButton{BeginFrame}{\begin{frame}}}%
%BeginExpansion
\begin{frame}%
%EndExpansion

\QTR{frametitle}{Caso Aula: China Evergrande 2021}

%TCIMACRO{\TeXButton{Transparency}{\setbeamercovered{transparent=20}}}%
%BeginExpansion
\setbeamercovered{transparent=20}%
%EndExpansion

Siamo nell'estate 2021 e osserviamo il problema dalla prospettiva di un risk manager.

\vspace{0.3cm}

Il portafoglio contiene:

\begin{itemize}

\item \(36\) debitori;

\item un investimento complessivo di \(200\) MUSD;

\item una posizione particolarmente concentrata su China Evergrande;

\item esposizioni distribuite tra diversi rating creditizi.

\end{itemize}

\vspace{0.4cm}

La domanda quantitativa \`e:

\begin{center}
\textbf{come cambia la distribuzione della perdita di portafoglio se
introduciamo un limite single-name del \(10\%\)?}
\end{center}

\vspace{0.3cm}

Il confronto deve tenere conto contemporaneamente di:

\[
\text{regime sistemico}
\;+\;
\text{migrazioni creditizie}
\;+\;
\text{default}
\;+\;
\text{severit\`a}.
\]

%TCIMACRO{\TeXButton{Transition: Box Out}{\transboxout}}}%
%BeginExpansion
\transboxout%
%EndExpansion
%TCIMACRO{\TeXButton{EndFrame}{\end{frame}}}%
%BeginExpansion
\end{frame}%
%EndExpansion


%******************
% Slide 13
%******************

%TCIMACRO{\TeXButton{BeginFrame}{\begin{frame}}}%
%BeginExpansion
\begin{frame}%
%EndExpansion

\QTR{frametitle}{Il portafoglio iniziale: dove nasce la concentrazione}

%TCIMACRO{\TeXButton{Transparency}{\setbeamercovered{transparent=20}}}%
%BeginExpansion
\setbeamercovered{transparent=20}%
%EndExpansion

Il portafoglio iniziale ha valore complessivo $V_0^P=200\text{ MUSD}$.

La composizione \`e:

\begin{center}
\begin{tabular}{lccc}
\hline
\textbf{Classe} & \textbf{Numero} & \textbf{MUSD per nome} & \textbf{Totale} \\
\hline
Evergrande, \(B\) & 1 & 30 & 30 \\
\(BBB\) & 10 & 6 & 60 \\
\(BB\) & 10 & 5 & 50 \\
altri \(B\) & 10 & 4 & 40 \\
\(CCC\) & 5 & 4 & 20 \\
\hline
 & 36 &  & 200 \\
\hline
\end{tabular}
\end{center}

\vspace{0.3cm}

La quota iniziale di Evergrande \`e

\[
\frac{30}{200}=15\%.
\]

\begin{center}
\textbf{Il problema di concentrazione nasce dalla dimensione della
posizione rispetto al portafoglio.}
\end{center}

%TCIMACRO{\TeXButton{Transition: Box Out}{\transboxout}}}%
%BeginExpansion
\transboxout%
%EndExpansion
%TCIMACRO{\TeXButton{EndFrame}{\end{frame}}}%
%BeginExpansion
\end{frame}%
%EndExpansion


%******************
% Slide 14
%******************

%TCIMACRO{\TeXButton{BeginFrame}{\begin{frame}}}%
%BeginExpansion
\begin{frame}%
%EndExpansion

\QTR{frametitle}{Il modello: due livelli stocastici}

%TCIMACRO{\TeXButton{Transparency}{\setbeamercovered{transparent=20}}}%
%BeginExpansion
\setbeamercovered{transparent=20}%
%EndExpansion

Il modello combina una dinamica sistemica e una dinamica creditizia.

\vspace{0.3cm}

\textbf{Livello 1: regime sistemico}

\[
M_t\in\{N,S,C\},
\qquad
M_0=S.
\]

Il regime evolve secondo la matrice di transizione \(Q\).

\vspace{0.4cm}

\textbf{Livello 2: rating dei debitori}

\[
X_{i,t}\in\{BBB,BB,B,CCC,D\},
\]

con \(D\) stato assorbente.

La matrice utilizzata nella transizione \(t\to t+1\) dipende da \(M_t\):

\[
M_t
\longrightarrow
P^{(M_t)}
\longrightarrow
X_{i,t+1}.
\]

\vspace{0.3cm}

\begin{center}
\textbf{Il regime comune modifica la dinamica creditizia di tutti i
debitori.}
\end{center}

%TCIMACRO{\TeXButton{Transition: Box Out}{\transboxout}}}%
%BeginExpansion
\transboxout%
%EndExpansion
%TCIMACRO{\TeXButton{EndFrame}{\end{frame}}}%
%BeginExpansion
\end{frame}%
%EndExpansion

%******************
% Slide 15
%******************

%TCIMACRO{\TeXButton{BeginFrame}{\begin{frame}}}%
%BeginExpansion
\begin{frame}%
%EndExpansion

\QTR{frametitle}{Default e severit\`a: il regime conta due volte}

%TCIMACRO{\TeXButton{Transparency}{\setbeamercovered{transparent=20}}}%
%BeginExpansion
\setbeamercovered{transparent=20}%
%EndExpansion

Definiamo il tempo di default del debitore \(i\):

\[
\tau_i
=
\min
\left\{
t\in\{1,\ldots,4\}:X_{i,t}=D
\right\}.
\]

\vspace{0.3cm}

Il regime sistemico influenza la perdita attraverso due canali.

\begin{enumerate}

\item \textbf{Probabilit\`a di migrazione e default}

La matrice di transizione del rating dipende da \(M_t\).

\item \textbf{Severit\`a della perdita}

La LGD dipende dal regime sistemico:

\[
\operatorname{LGD}^{(N)}=0.55,
\qquad
\operatorname{LGD}^{(S)}=0.70,
\qquad
\operatorname{LGD}^{(C)}=0.85.
\]

\end{enumerate}

\vspace{0.3cm}

Se il debitore entra in default al tempo \(\tau_i\), il default si
realizza nella transizione da \(\tau_i-1\) a \(\tau_i\).

La LGD applicata \`e, quindi, quella associata al regime sistemico

\[
M_{\tau_i-1}.
\]

%TCIMACRO{\TeXButton{Transition: Box Out}{\transboxout}}}%
%BeginExpansion
\transboxout%
%EndExpansion
%TCIMACRO{\TeXButton{EndFrame}{\end{frame}}}%
%BeginExpansion
\end{frame}%
%EndExpansion

%******************
% Slide 16
%******************

%TCIMACRO{\TeXButton{BeginFrame}{\begin{frame}}}%
%BeginExpansion
\begin{frame}%
%EndExpansion

\QTR{frametitle}{Dalla traiettoria creditizia alla perdita}

%TCIMACRO{\TeXButton{Transparency}{\setbeamercovered{transparent=20}}}%
%BeginExpansion
\setbeamercovered{transparent=20}%
%EndExpansion

Per ogni debitore \(i\):

\begin{itemize}

\item \(V_{i,0}\): esposizione iniziale, misurata in MUSD;

\item \(\ell(X_{i,0},X_{i,4})\): quota dell'esposizione persa per
effetto della migrazione di rating, quindi una quantit\`a adimensionale;

\item \(EAD_i\): esposizione al momento del default, misurata in MUSD;

\item \(\operatorname{LGD}\): quota dell'EAD persa in caso di default,
quindi una quantit\`a adimensionale.

\end{itemize}

%\vspace{0.2cm}

Nel caso didattico, assumiamo $EAD_i=V_{i,0}$.

Se non si verifica default:

\[
L_i
=
V_{i,0}\,
\ell(X_{i,0},X_{i,4}),
\qquad
\tau_i>4.
\]

Se si verifica default:

\[
L_i
=
V_{i,0}\,
\operatorname{LGD}^{(M_{\tau_i-1})},
\qquad
\tau_i\leq 4.
\]

%\vspace{0.2cm}

%In entrambi i casi:

%\[
%\text{MUSD}\times\text{quota di perdita}
%=
%\text{MUSD}.
%\]

La perdita di portafoglio \`e, quindi, $L^P=\sum_{i=1}^{36}L_i$.

%TCIMACRO{\TeXButton{Transition: Box Out}{\transboxout}}}%
%BeginExpansion
\transboxout%
%EndExpansion
%TCIMACRO{\TeXButton{EndFrame}{\end{frame}}}%
%BeginExpansion
\end{frame}%
%EndExpansion


%******************
% Slide 17
%******************

%TCIMACRO{\TeXButton{BeginFrame}{\begin{frame}}}%
%BeginExpansion
\begin{frame}%
%EndExpansion

\QTR{frametitle}{La politica di concentrazione: dal 15\% al 10\%}

%TCIMACRO{\TeXButton{Transparency}{\setbeamercovered{transparent=20}}}%
%BeginExpansion
\setbeamercovered{transparent=20}%
%EndExpansion

Imponiamo il vincolo

\[
\max_i
\frac{V_{i,0}}{V_0^P}
\leq 10\%.
\]

%\vspace{0.3cm}

Poich\'e il portafoglio vale \(200\) MUSD, la posizione massima
ammessa \`e

\[
20\text{ MUSD}.
\]

Evergrande passa quindi da

\[
30\text{ MUSD}
\qquad\text{a}\qquad
20\text{ MUSD}.
\]

%\vspace{0.3cm}

I \(10\) MUSD liberati vengono riallocati:

\begin{itemize}

\item esclusivamente ai debitori inizialmente \(BBB\) e \(BB\);

\item in proporzione agli investimenti iniziali nelle due classi.

\end{itemize}

%\vspace{0.3cm}

Il capitale complessivo resta invariato:

\[
V_0^{P,\mathrm{base}}
=
V_0^{P,\mathrm{lim}}
=
200\text{ MUSD}.
\]

%TCIMACRO{\TeXButton{Transition: Box Out}{\transboxout}}}%
%BeginExpansion
\transboxout%
%EndExpansion
%TCIMACRO{\TeXButton{EndFrame}{\end{frame}}}%
%BeginExpansion
\end{frame}%
%EndExpansion


%******************
% Slide 18
%******************

%TCIMACRO{\TeXButton{BeginFrame}{\begin{frame}}}%
%BeginExpansion
\begin{frame}%
%EndExpansion

\QTR{frametitle}{Confrontare i due portafogli sugli stessi scenari}

%TCIMACRO{\TeXButton{Transparency}{\setbeamercovered{transparent=20}}}%
%BeginExpansion
\setbeamercovered{transparent=20}%
%EndExpansion

Per confrontare correttamente il portafoglio base e quello soggetto al
limite, utilizziamo gli stessi scenari simulati.

\vspace{0.3cm}

In ogni replica \(r\):

\begin{itemize}

\item la traiettoria sistemica \`e la stessa;

\item le traiettorie creditizie dei debitori sono le stesse;

\item cambia soltanto il vettore degli investimenti.

\end{itemize}

\vspace{0.4cm}

Possiamo quindi definire, scenario per scenario,

\[
\Delta L^{(r)}
=
L_{\mathrm{lim}}^{P,(r)}
-
L_{\mathrm{base}}^{P,(r)}.
\]

\vspace{0.3cm}

Questa scelta consente di isolare l'effetto della politica di
concentrazione dagli effetti della variabilit\`a casuale della
simulazione.

\vspace{0.3cm}

\begin{center}
\textbf{Stessi shock, portafogli diversi: il confronto diventa
economicamente interpretabile.}
\end{center}

%TCIMACRO{\TeXButton{Transition: Box Out}{\transboxout}}}%
%BeginExpansion
\transboxout%
%EndExpansion
%TCIMACRO{\TeXButton{EndFrame}{\end{frame}}}%
%BeginExpansion
\end{frame}%
%EndExpansion


%============================================================
% SEZIONE 4 — CASO AULA: DALLA STRUTTURA TEORICA AL NOTEBOOK
%============================================================

\section{Caso Aula: dalla struttura teorica al notebook}

%******************
% Slide 19
%******************

%TCIMACRO{\TeXButton{BeginFrame}{\begin{frame}}}%
%BeginExpansion
\begin{frame}%
%EndExpansion

\QTR{frametitle}{Dal modello alla sequenza operativa}

%TCIMACRO{\TeXButton{Transparency}{\setbeamercovered{transparent=20}}}%
%BeginExpansion
\setbeamercovered{transparent=20}%
%EndExpansion

La struttura teorica del caso deve ora diventare una sequenza
computazionale.

\vspace{0.3cm}

\[
\text{regime sistemico}
\longrightarrow
\text{migrazioni creditizie}
\longrightarrow
\text{tempi di default}
\]

\[
\longrightarrow
\text{perdite individuali}
\longrightarrow
\text{perdita di portafoglio}
\longrightarrow
\text{misure di rischio}.
\]

\vspace{0.4cm}

Nel notebook dovremo:

\begin{enumerate}

\item simulare le traiettorie del regime \(M_t\);

\item simulare le traiettorie creditizie \(X_{i,t}\);

\item individuare gli eventuali tempi di default \(\tau_i\);

\item trasformare ogni traiettoria in una perdita monetaria;

\item aggregare le perdite e confrontare i due portafogli.

\end{enumerate}

\vspace{0.3cm}

\begin{center}
\textbf{Il codice deve rispettare l'ordine logico del modello.}
\end{center}

%TCIMACRO{\TeXButton{Transition: Box Out}{\transboxout}}}%
%BeginExpansion
\transboxout%
%EndExpansion
%TCIMACRO{\TeXButton{EndFrame}{\end{frame}}}%
%BeginExpansion
\end{frame}%
%EndExpansion


%******************
% Slide 20
%******************

%TCIMACRO{\TeXButton{BeginFrame}{\begin{frame}}}%
%BeginExpansion
\begin{frame}%
%EndExpansion

\QTR{frametitle}{Quali output dovremo ottenere?}

%TCIMACRO{\TeXButton{Transparency}{\setbeamercovered{transparent=20}}}%
%BeginExpansion
\setbeamercovered{transparent=20}%
%EndExpansion

La simulazione Monte Carlo produce due distribuzioni di perdita:

\[
L_{\mathrm{base}}^P
\qquad\text{e}\qquad
L_{\mathrm{lim}}^P.
\]

\vspace{0.3cm}

Per entrambe, dovremo calcolare:

\begin{itemize}

\item perdita media, varianza e deviazione standard;

\item VaR e CVaR ai livelli di confidenza previsti;

\item perdita media negli scenari con crisi e senza crisi.

\end{itemize}

\vspace{0.3cm}

Il confronto tra le due politiche richiede inoltre:

\[
\Delta L^{(r)}
=
L_{\mathrm{lim}}^{P,(r)}
-
L_{\mathrm{base}}^{P,(r)}.
\]

\vspace{0.3cm}

\begin{center}
\textbf{Non cerchiamo un solo numero: confrontiamo due distribuzioni
di perdita e il loro comportamento in coda.}
\end{center}

%TCIMACRO{\TeXButton{Transition: Box Out}{\transboxout}}}%
%BeginExpansion
\transboxout%
%EndExpansion
%TCIMACRO{\TeXButton{EndFrame}{\end{frame}}}%
%BeginExpansion
\end{frame}%
%EndExpansion


%******************
% Slide 21
%******************

%TCIMACRO{\TeXButton{BeginFrame}{\begin{frame}}}%
%BeginExpansion
\begin{frame}%
%EndExpansion

\QTR{frametitle}{Prima degli output: che cosa dobbiamo controllare?}

%TCIMACRO{\TeXButton{Transparency}{\setbeamercovered{transparent=20}}}%
%BeginExpansion
\setbeamercovered{transparent=20}%
%EndExpansion

Un notebook corretto non deve soltanto produrre risultati: deve
verificare che la simulazione rispetti il modello.

%\vspace{0.3cm}

I controlli principali riguardano:

\begin{itemize}

\item correttezza delle matrici di transizione e dello stato iniziale;

\item assorbimento dello stato \(D\);

\item coerenza tra \(\tau_i\), regime sistemico e LGD applicata;

\item conservazione del valore complessivo del portafoglio dopo
l'applicazione del limite;

\item utilizzo degli stessi scenari per portafoglio base e portafoglio
con limite;

\item coerenza dimensionale delle quantit\`a monetarie e delle misure di
rischio.

\end{itemize}

%\vspace{0.4cm}

\begin{center}
\textbf{Da questo punto il lavoro prosegue direttamente in
VS Code/Jupyter.}
\end{center}

%\vspace{0.2cm}

Il notebook dovr\`a costruire progressivamente simulazioni, output,
controlli e interpretazione finale.

%TCIMACRO{\TeXButton{Transition: Box Out}{\transboxout}}}%
%BeginExpansion
\transboxout%
%EndExpansion
%TCIMACRO{\TeXButton{EndFrame}{\end{frame}}}%
%BeginExpansion
\end{frame}%
%EndExpansion

%============================================================
% SEZIONE 5 — CASO TAKEHOME: LEHMAN BROTHERS 2008
%============================================================

\section{Caso TakeHome: Lehman Brothers 2008}

%******************
% Slide 22
%******************

%TCIMACRO{\TeXButton{BeginFrame}{\begin{frame}}}%
%BeginExpansion
\begin{frame}%
%EndExpansion

\QTR{frametitle}{Caso TakeHome: Lehman Brothers 2008}

%TCIMACRO{\TeXButton{Transparency}{\setbeamercovered{transparent=20}}}%
%BeginExpansion
\setbeamercovered{transparent=20}%
%EndExpansion

Siamo nel giugno 2008 e analizziamo un portafoglio creditizio con una
posizione rilevante su Lehman Brothers.

%\vspace{0.3cm}

Il portafoglio iniziale contiene:

\begin{itemize}

\item \(40\) debitori;

\item un investimento complessivo di \(173\) MUSD;

\item una posizione Lehman di \(25\) MUSD;

\item esposizioni distribuite tra rating \(A\), \(BBB\), \(BB\) e \(B\).

\end{itemize}

%\vspace{0.4cm}

Il controfattuale elimina la posizione Lehman e redistribuisce il
capitale sui \(39\) debitori residui, mantenendo invariato il valore
complessivo del portafoglio.

%\vspace{0.3cm}

La domanda quantitativa \`e:

\begin{center}
\textbf{come cambia la distribuzione della perdita quando la
concentrazione su Lehman viene eliminata?}
\end{center}

%\vspace{0.3cm}

Anche in questo caso il confronto dovr\`a riguardare l'intera
distribuzione delle perdite, non soltanto la perdita media.

%TCIMACRO{\TeXButton{Transition: Box Out}{\transboxout}}}%
%BeginExpansion
\transboxout%
%EndExpansion
%TCIMACRO{\TeXButton{EndFrame}{\end{frame}}}%
%BeginExpansion
\end{frame}%
%EndExpansion


%******************
% Slide 23
%******************

%TCIMACRO{\TeXButton{BeginFrame}{\begin{frame}}}%
%BeginExpansion
\begin{frame}%
%EndExpansion

\QTR{frametitle}{Che cosa resta uguale e che cosa cambia?}

%TCIMACRO{\TeXButton{Transparency}{\setbeamercovered{transparent=20}}}%
%BeginExpansion
\setbeamercovered{transparent=20}%
%EndExpansion

L'architettura metodologica resta la stessa del Caso Aula:

\begin{itemize}

\item regime sistemico markoviano;

\item migrazioni creditizie condizionate al regime;

\item simulazione Monte Carlo;

\item costruzione della perdita di portafoglio;

\item VaR, CVaR e confronto controfattuale;

\item stessi scenari per i debitori comuni ai due portafogli.

\end{itemize}

%\vspace{0.3cm}

Cambiano invece alcuni elementi del modello:

\begin{itemize}

\item il regime sistemico assume gli stati \(O,S,C\);

\item la LGD \`e costante e pari a \(0.65\);

\item la posizione Lehman viene eliminata completamente;

\item il capitale viene redistribuito proporzionalmente sui \(39\)
debitori residui.

\end{itemize}

%\vspace{0.3cm}

\begin{center}
\textbf{Il TakeHome riutilizza lo stesso metodo, ma non riproduce lo
stesso problema.}
\end{center}

%\vspace{0.2cm}

La riallocazione modifica anche la composizione per rating del
portafoglio: il confronto non \`e quindi un puro esperimento sulla
concentrazione.

%TCIMACRO{\TeXButton{Transition: Box Out}{\transboxout}}}%
%BeginExpansion
\transboxout%
%EndExpansion
%TCIMACRO{\TeXButton{EndFrame}{\end{frame}}}%
%BeginExpansion
\end{frame}%
%EndExpansion

\end{document}